\documentclass[10pt,a4paper]{amsart}
\usepackage[margin=19mm]{geometry}
\usepackage{amsmath,amssymb,mathtools}
\usepackage[scr=rsfs,cal=euler]{mathalfa}
\usepackage{xfrac}
\usepackage[unicode,hidelinks,bookmarks=false]{hyperref}
\newcommand{\ie}{i.e.\ }
\newcommand{\cf}{cf.\ }
\newcommand{\resp}{respectively\ }
\newcommand{\Subsection}{\subsection}
\providecommand{\nobreakdash}{\nobreak\hbox{-}}
\setlength{\emergencystretch}{2em}
\multlinegap0pt
\usepackage{bm}
\usepackage{bbm}
\usepackage{dsfont}
\usepackage{xspace}
\usepackage{cancel}
\usepackage{mathtools}
\usepackage{amsthm}
\makeatletter
\@ifundefined{theorem}{\newtheorem{theorem}{Theorem}}{}
\@ifundefined{definition}{\newtheorem{definition}[theorem]{Definition}}{}
\@ifundefined{proposition}{\newtheorem{proposition}[theorem]{Proposition}}{}
\makeatother
\newcommand{\Comm}{\mathrm{Comm}}
\newcommand{\Z}{\mathrm{Z}}
\title[Higher Commutativity in Finite Groups, Rigidity, Extremal bo]{Higher Commutativity in Finite Groups, Rigidity, Extremal bounds, and Heisenberg-Type Families}
\author{Vadim E Levit, Robert Shwartz}
\date{Source: arXiv:2605.17171v1 (CC BY 4.0)}
\begin{document}
\maketitle

\noindent\emph{Source and licence.} Higher Commutativity in Finite Groups, Rigidity, Extremal bounds, and Heisenberg-Type Families. Version arXiv:2605.17171v1 (CC BY 4.0), distributed under the Creative Commons Attribution 4.0 International licence, as recorded in the arXiv licence metadata for this exact version and shown on the abstract page https://arxiv.org/abs/2605.17171v1. Full rights notice: licenses/arxiv-2605.17171-CC-BY-4.0.txt.\par\smallskip

\noindent\textit{Editorial scope.} Counted unit: the Theorem whose source label is \texttt{thm:heisenberg-full-hierarchy} in the article (source0.tex lines 2010--2073, source label \texttt{thm:heisenberg-full-hierarchy}), delivered together with the article's own complete original proof of it (source0.tex lines 2075--2302, one \texttt{proof} environment of 228 lines). No mathematics is written, repaired or re-derived: statement and proof are the article's own text line for line. Required context retained uncounted: def:class2-exp-p (lines 1853--1869); prop:class2-reduction (lines 1871--1885); def:heisenberg-type (lines 1953--1967). Every definition, display and label of those lines is retained unchanged. Omitted from this package and not redistributed: 1--1852, 1870--1870, 1886--1952, 1968--2009, 2074--2074, 2303--2624. Nothing inside a retained excerpt is omitted or altered: every retained body is delivered exactly as the article's lines are written, so no omission receipt of any kind accompanies this package. Citations: the retained excerpts carry no citation command at all, so no bibliography entry of the article is transcribed and no reference is redistributed. Reference adaptation: none was needed and none was made; every reference inside the delivered unit resolves inside it, so no unresolved reference or citation is produced. Printed slips kept literal: no printed word, symbol, hypothesis or number of the counted statement or of its proof was altered. Layout and class: the article's own class is \texttt{article} with packages including amsmath, amssymb, amsthm, mathrsfs, geometry, hyperref; the wrapper is the lane's pinned \texttt{amsart} editorial wrapper at 10pt on A4, which restates the theorem-like environments the retained text needs and adds the article's own macro definitions that the retained text needs (\texttt{\textbackslash Comm}, \texttt{\textbackslash Z}), with extra wrapper packages (bm, bbm, dsfont, xspace, cancel, mathtools). The delivered numbering is the unit's own. Assets: no retained span contains a figure environment, a graphics-inclusion command, a PDF-page inclusion, a table or a TikZ picture; no image file is copied into this package, the package declares no asset dependency and no image file is redistributed with it. Licence: version arXiv:2605.17171v1 of this article is distributed under the Creative Commons Attribution 4.0 International licence, as recorded in the arXiv licence metadata for this exact version and shown on the abstract page https://arxiv.org/abs/2605.17171v1. A full rights notice is delivered at licenses/arxiv-2605.17171-CC-BY-4.0.txt. Characters: every retained excerpt is pure ASCII, so the delivered engine, XeLaTeX with its default Latin Modern text font, reports no missing character.
\begin{definition}[Class-$2$ exponent-$p$ groups and the commutator map]\label{def:class2-exp-p}
Let $p$ be a prime and let $G$ be a finite $p$-group of nilpotency class $2$ and exponent $p$.
(If $p=2$ then exponent $p$ forces $G$ abelian, so the non-abelian case is relevant only for odd $p$.)
Set
\[
V:=G/\Z(G),\qquad W:=G'.
\]
Then $V$ and $W$ are naturally $\mathbb F_p$-vector spaces, and the commutator induces an alternating $\mathbb F_p$-bilinear map
\[
\beta:V\times V\to W,\qquad \beta(x\Z(G),y\Z(G))=[x,y],
\]
equivalently a linear map of $\mathbb F_p$-vector spaces
\[
\widetilde\beta:\Lambda^2 V\to W.
\]
We call $\widetilde\beta$ (or $\beta$) the \emph{commutator tensor} of $G$.
\end{definition}

\begin{proposition}[Reduction to isotropic tuples]\label{prop:class2-reduction}
Let $G$ be a finite $p$-group of class $2$ and exponent $p$, with commutator tensor $\widetilde\beta:\Lambda^2V\to W$
as in Definition~\ref{def:class2-exp-p}.
For $r\ge 1$ set
\[
N_r(V,\beta):=\big|\{(v_1,\dots,v_r)\in V^r:\ \beta(v_i,v_j)=0\ \forall i<j\}\big|.
\]
Then
\[
|\Comm_r(G)|=|\Z(G)|^r\,N_r(V,\beta),
\qquad\text{and hence}\qquad
P_r(G)=\frac{N_r(V,\beta)}{|V|^r}.
\]
In particular, $P_r(G)$ depends only on the alternating bilinear map $\beta$ on $V$ (and not on $|\Z(G)|$).
\end{proposition}

\begin{definition}[$\mathbb F_{p^m}$-Heisenberg type]\label{def:heisenberg-type}
Let $p$ be an odd prime and let $q:=p^m$.
We say that a finite $p$-group $G$ is of \emph{$\mathbb F_q$-Heisenberg type of rank $n$} if:
\begin{enumerate}
\item $G$ has nilpotency class $2$ and exponent $p$;
\item $G'=\Z(G)$ is elementary abelian of order $q$ (so $G'\cong \mathbb F_q$ as additive groups);
\item $V:=G/\Z(G)$ carries the structure of an $\mathbb F_q$-vector space of dimension $2n$ such that,
under an identification $G'\cong \mathbb F_q$, the commutator map
\(
\beta:V\times V\to \mathbb F_q
\)
is a nondegenerate alternating $\mathbb F_q$-bilinear form.
\end{enumerate}
Equivalently, at the level of commutator geometry, $G$ is represented by a nondegenerate alternating $\mathbb F_q$-bilinear form on $\mathbb F_q^{2n}$; this determines the associated isoclinism data, not a preferred group isomorphism type.
\end{definition}

\begin{theorem}[Closed hierarchy and rigidity in the $\mathbb F_q$-Heisenberg family]\label{thm:heisenberg-full-hierarchy}\label{thm:rank2-heisenberg-isoclinism}
Let $p$ be an odd prime, let $q:=p^m$, and let $G$ be of $\mathbb F_q$-Heisenberg type of rank $n$
in the sense of Definition~\ref{def:heisenberg-type}. Write
\[
V:=G/\Z(G)\cong \mathbb F_q^{2n}.
\]
For $0\le k\le n$, let $L_{n,k}(q)$ denote the number of $k$-dimensional totally isotropic
$\mathbb F_q$-subspaces of $V$. Then
\[
L_{n,k}(q)=\prod_{i=0}^{k-1}\frac{q^{2n-2i}-1}{q^{k-i}-1}.
\]

For every $r\ge 1$ one has
\[
I_r^{(q)}(n)
=
\sum_{k=0}^{\min(n,r)}
L_{n,k}(q)\prod_{i=0}^{k-1}(q^r-q^i),
\]
and hence
\[
|\Comm_r(G)|
=
|\Z(G)|^r
\sum_{k=0}^{\min(n,r)}
L_{n,k}(q)\prod_{i=0}^{k-1}(q^r-q^i),
\]
equivalently
\[
P_r(G)
=
q^{-2nr}
\sum_{k=0}^{\min(n,r)}
L_{n,k}(q)\prod_{i=0}^{k-1}(q^r-q^i).
\]
In particular, the full hierarchy $\{P_r(G)\}_{r\ge2}$ depends only on $(q,n)$.

Now let $q':=p^{m'}$, and let $H$ be another finite $p$-group of $\mathbb F_{q'}$-Heisenberg type of rank $n'$.
Then the following are equivalent:
\begin{enumerate}
\item $G$ and $H$ are isoclinic;
\item $(q,n)=(q',n')$;
\item $P_2(G)=P_2(H)$ and $P_3(G)=P_3(H)$;
\item $P_r(G)=P_r(H)$ for all $r\ge2$.
\end{enumerate}
Moreover, in (3) the parameter $q$ is the unique positive root of
\[
\bigl(P_2(G)-P_3(G)\bigr)X^2+\bigl(P_2(G)-1\bigr)X+\bigl(P_2(G)-1\bigr)=0,
\]
and then
\[
q^{-2n}=\frac{qP_2(G)-1}{q-1}.
\]

Finally, with the manuscript convention
\[
\mathcal P_G(z):=\sum_{r\ge2}P_r(G)z^{r-2},
\]
the series $\mathcal P_G(z)$ is rational and its pole set is contained in
\[
\{q^n,q^{n+1},\dots,q^{2n}\}.
\]
In particular, the pole at $z=q^n$ is present.
\end{theorem}

\begin{proof}
Let $\langle\cdot,\cdot\rangle$ be the nondegenerate alternating $\mathbb F_q$-bilinear form on
$V\cong \mathbb F_q^{2n}$ coming from the commutator map.

First we compute $L_{n,k}(q)$.
Let $\mathcal B_{n,k}(q)$ be the set of ordered $k$-tuples
\[
(u_1,\dots,u_k)\in V^k
\]
such that $u_1,\dots,u_k$ are linearly independent and span a totally isotropic subspace.
We count $\mathcal B_{n,k}(q)$ in two ways.

Choose such a tuple inductively.
After choosing $u_1,\dots,u_{j-1}$, their span
\[
U_{j-1}:=\langle u_1,\dots,u_{j-1}\rangle
\]
is totally isotropic of dimension $j-1$, hence $U_{j-1}\subseteq U_{j-1}^{\perp}$ and
\[
\dim_{\mathbb F_q}(U_{j-1}^{\perp})=2n-(j-1).
\]
To keep the enlarged span totally isotropic and linearly independent, the next vector must lie in
$U_{j-1}^{\perp}\setminus U_{j-1}$. Therefore the number of choices for $u_j$ is
\[
|U_{j-1}^{\perp}|-|U_{j-1}|=q^{2n-j+1}-q^{j-1}.
\]
Thus
\[
|\mathcal B_{n,k}(q)|
=
\prod_{j=1}^{k}\bigl(q^{2n-j+1}-q^{j-1}\bigr).
\]

On the other hand, each $k$-dimensional totally isotropic subspace has exactly
\[
\prod_{j=1}^{k}(q^k-q^{j-1})
\]
ordered bases. Hence
\[
L_{n,k}(q)
=
\frac{\prod_{j=1}^{k}(q^{2n-j+1}-q^{j-1})}
     {\prod_{j=1}^{k}(q^k-q^{j-1})}.
\]
Factoring out $q^{j-1}$ from numerator and denominator in each factor gives
\[
L_{n,k}(q)
=
\prod_{j=1}^{k}\frac{q^{2n-2j+2}-1}{q^{k-j+1}-1}
=
\prod_{i=0}^{k-1}\frac{q^{2n-2i}-1}{q^{k-i}-1},
\]
as claimed.

Now let
\[
N_r(V):=\bigl|\{(v_1,\dots,v_r)\in V^r:\ \langle v_i,v_j\rangle=0\ \text{for all }i<j\}\bigr|.
\]
By definition, $N_r(V)=I_r^{(q)}(n)$.
We count $N_r(V)$ by the dimension of the span of the tuple.

If $(v_1,\dots,v_r)$ is pairwise orthogonal, then
\[
U:=\langle v_1,\dots,v_r\rangle
\]
is a totally isotropic subspace of $V$.
Let $k:=\dim U$. Then necessarily $0\le k\le \min(n,r)$.
Conversely, fix a $k$-dimensional totally isotropic subspace $U\le V$.
The pairwise orthogonal $r$-tuples whose span is exactly $U$ are precisely the ordered $r$-tuples in $U^r$
that generate $U$.

Choose an $\mathbb F_q$-linear isomorphism $U\cong \mathbb F_q^k$.
Then ordered generating $r$-tuples of $U$ are in bijection with surjective linear maps
\[
T:\mathbb F_q^r\twoheadrightarrow \mathbb F_q^k,
\qquad
T(e_i)=v_i.
\]
The number of such surjections equals the number of rank-$k$ $k\times r$ matrices over $\mathbb F_q$.
Transposing, this is the number of injective linear maps
\[
\mathbb F_q^k\hookrightarrow \mathbb F_q^r,
\]
which is
\[
\prod_{i=0}^{k-1}(q^r-q^i):
\]
choose the image of the first basis vector in $q^r-1$ ways, the second outside its span in $q^r-q$ ways,
and so on.

Therefore, for each fixed $k$-dimensional totally isotropic subspace $U$, the number of pairwise orthogonal
$r$-tuples spanning $U$ is
\[
\prod_{i=0}^{k-1}(q^r-q^i).
\]
Summing over all such $U$ gives
\[
N_r(V)
=
\sum_{k=0}^{\min(n,r)}
L_{n,k}(q)\prod_{i=0}^{k-1}(q^r-q^i).
\]
Since $N_r(V)=I_r^{(q)}(n)$, this proves the formula for $I_r^{(q)}(n)$.

Now apply Proposition~\ref{prop:class2-reduction}. Because $|V|=q^{2n}$,
\[
|\Comm_r(G)|=|\Z(G)|^r\,N_r(V)
\]
and
\[
P_r(G)=\frac{N_r(V)}{|V|^r}=q^{-2nr}N_r(V),
\]
which yields the displayed formulas for $|\Comm_r(G)|$ and $P_r(G)$.
In particular, $P_r(G)$ depends only on $(q,n)$.

We now prove the rigidity statement.

Assume first that $G$ and $H$ are isoclinic.
Then $G'\cong H'$ as groups, so
\[
|G'|=|H'|,
\]
hence $q=q'$.
Also
\[
G/\Z(G)\cong H/\Z(H),
\]
so
\[
q^{2n}=|G:\Z(G)|=|H:\Z(H)|=(q')^{2n'}=q^{2n'}.
\]
Therefore $n=n'$, proving $(1)\Rightarrow(2)$.

Conversely, assume $(q,n)=(q',n')$.
Choose identifications
\[
G'\cong \mathbb F_q,\qquad H'\cong \mathbb F_q.
\]
Then $V_G:=G/\Z(G)$ and $V_H:=H/\Z(H)$ are both $2n$-dimensional symplectic spaces over $\mathbb F_q$.
Choose symplectic bases
\[
(e_1,\dots,e_n,f_1,\dots,f_n)\quad\text{for }V_G,
\]
\[
(e_1',\dots,e_n',f_1',\dots,f_n')\quad\text{for }V_H.
\]
The unique $\mathbb F_q$-linear map $\phi:V_G\to V_H$ sending $e_i\mapsto e_i'$ and $f_i\mapsto f_i'$
is an isometry of alternating forms.
Let $\psi:G'\to H'$ be the chosen identification.
Then
\[
\psi\circ \beta_G=\beta_H\circ(\phi\times\phi),
\]
so $G$ and $H$ are isoclinic. Thus $(2)\Rightarrow(1)$.

The implication $(2)\Rightarrow(4)$ is immediate from the closed formula for $P_r$,
and $(4)\Rightarrow(3)$ is trivial.

It remains to prove $(3)\Rightarrow(2)$.
Taking $r=2$ and $r=3$ in the closed formula yields
\[
P_2(G)=\frac1q+\Bigl(1-\frac1q\Bigr)q^{-2n},
\qquad
P_3(G)=\frac1{q^3}+\Bigl(1-\frac1{q^3}\Bigr)q^{-2n}.
\]
Set
\[
a:=q^{-2n}.
\]
Then
\[
P_2(G)=q^{-1}+(1-q^{-1})a,\qquad
P_3(G)=q^{-3}+(1-q^{-3})a.
\]
Subtracting and comparing with $1-P_2(G)$ gives
\[
\frac{P_2(G)-P_3(G)}{1-P_2(G)}
=
\frac{q^{-1}-q^{-3}}{1-q^{-1}}
=
q^{-1}+q^{-2}.
\]
Thus $q$ is determined uniquely by $P_2(G)$ and $P_3(G)$, because the function
$x\mapsto x^{-1}+x^{-2}$ is strictly decreasing on $(0,\infty)$.
Equivalently, eliminating $a$ yields the quadratic
\[
\bigl(P_2(G)-P_3(G)\bigr)X^2+\bigl(P_2(G)-1\bigr)X+\bigl(P_2(G)-1\bigr)=0,
\]
whose unique positive root is $X=q$.

Once $q$ is known, the formula for $P_2(G)$ gives
\[
a=q^{-2n}=\frac{qP_2(G)-1}{q-1}.
\]
Since the map $n\mapsto q^{-2n}$ is injective on $\mathbb Z_{\ge0}$, this determines $n$.
Hence equality of $P_2$ and $P_3$ for $G$ and $H$ forces $(q,n)=(q',n')$, proving $(3)\Rightarrow(2)$.
This completes the equivalence of (1)--(4).

Finally, for each fixed $k$ the factor
\[
\prod_{i=0}^{k-1}(q^r-q^i)
\]
is a polynomial in $q^r$ of degree $k$.
Therefore there exist coefficients $c_0,\dots,c_n\in \mathbb Q(q)$, depending only on $(q,n)$, such that
\[
P_r(G)=\sum_{j=0}^{n} c_j\,q^{-(2n-j)r}.
\]
Hence
\[
\mathcal P_G(z)
=
\sum_{r\ge2}P_r(G)z^{r-2}
=
\sum_{j=0}^{n} c_j\sum_{r\ge2}q^{-(2n-j)r}z^{r-2}
=
\sum_{j=0}^{n}\frac{c_j\,q^{-2(2n-j)}}{1-z/q^{\,2n-j}}.
\]
Thus $\mathcal P_G(z)$ is rational and its poles lie among
\[
q^n,q^{n+1},\dots,q^{2n}.
\]
Moreover, the term $k=n$ contributes
\[
L_{n,n}(q)\,q^{-2nr}\cdot q^{nr}=L_{n,n}(q)\,q^{-nr},
\]
and no term with $k<n$ produces $q^{-nr}$.
So the coefficient of $q^{-nr}$ is $L_{n,n}(q)>0$, and the pole at $z=q^n$ is present.
\end{proof}

\end{document}
